\section{NỘI DUNG VÀ PHƯƠNG PHÁP}

\subsection{Quy trình triển khai}

Đề tài được triển khai theo quy trình gồm bốn lớp. Lớp thứ nhất là môi trường mạng giả lập dựa trên Kathará~\citep{kathara2019}, trong đó mỗi router, host hoặc switch được chạy trong container và có thể được tiêm lỗi có kiểm soát. Lớp thứ hai là hệ thống MCP server~\citep{mcp2024}, chuẩn hóa các thao tác chẩn đoán thành công cụ để agent gọi trong quá trình suy luận. Lớp thứ ba là workflow agent xây dựng bằng LangGraph~\citep{langgraph2024}, hỗ trợ các luồng ReAct~\citep{react2023}, Plan \& Execute và Reflexion~\citep{reflexion2023}. Lớp cuối cùng là hai mô-đun tự tiến hóa: Procedural Memory và Tool Refinement.

Một episode chẩn đoán diễn ra theo trình tự: khởi động lab mạng, tiêm lỗi, chạy agent, thu thập trace gọi công cụ, nộp kết quả gồm \code{is\_anomaly}, \code{faulty\_devices}, \code{root\_cause\_name}, sau đó đánh giá và kích hoạt học ngoại tuyến. Với baseline, mỗi episode độc lập. Với evolved agent, bộ nhớ kỹ năng và tài liệu công cụ đã được tinh chỉnh từ các episode trước được nạp lại vào ngữ cảnh để hỗ trợ case tiếp theo.

Luồng xử lý được tách thành hai pha có chủ đích. Trong \textit{Diagnosis Phase}, agent được quyền dùng đầy đủ công cụ để thu thập bằng chứng, kiểm tra giả thuyết và loại trừ nguyên nhân sai. Trong \textit{Submission Phase}, agent chỉ còn công cụ nộp kết quả, nhờ đó giảm rủi ro vừa chẩn đoán vừa thay đổi kết luận cuối cùng một cách thiếu kiểm soát. Thiết kế này bám sát tinh thần ReAct, nơi suy luận và hành động công cụ được xen kẽ để giải quyết tác vụ nhiều bước~\citep{react2023}, nhưng bổ sung lớp đánh giá sau episode để phục vụ học liên tục.

\subsection{Công cụ và công nghệ sử dụng}

Hệ thống sử dụng NIKA làm benchmark chẩn đoán mạng, Kathará cho ảo hóa topology, FastMCP cho server công cụ, LangGraph/LangChain cho điều phối agent, InfluxDB cho telemetry và Python 3.12 với \code{uv} để quản lý môi trường. Với các kịch bản định tuyến, hệ thống tương tác với FRRouting để đọc cấu hình và trạng thái route~\citep{frr2021}; với kịch bản data-plane lập trình được, hệ thống đọc log và bảng hành vi của P4/BMv2, dựa trên mô hình lập trình packet processor của P4~\citep{p4lang2014}. Các MCP server chính gồm:
\begin{itemize}
  \item \textbf{Kathará Base}: kiểm tra reachability, ping, iperf, cấu hình host, netstat và thực thi shell trong container.
  \item \textbf{FRR}: đọc cấu hình OSPF/BGP, bảng định tuyến và thực thi lệnh \code{vtysh}.
  \item \textbf{BMv2/P4}: đọc log switch, dump bảng luồng, counter và register.
  \item \textbf{Telemetry}: truy vấn InfluxDB để lấy chỉ số mạng theo thời gian.
  \item \textbf{Task Management}: nhận kết quả cuối cùng để chấm điểm tự động.
  \item \textbf{Tool Evolution}: cung cấp tài liệu công cụ đã được tinh chỉnh cho agent.
\end{itemize}

\subsection{Mô-đun Procedural Memory}

Procedural Memory được thiết kế để lưu trữ kinh nghiệm chẩn đoán lâu dài dưới dạng \code{ProceduralSkill}, lấy cảm hứng từ hướng procedural memory trong Skill-Pro~\citep{skillpro2024}. Mỗi kỹ năng gồm điều kiện kích hoạt, chuỗi bước thực thi, điều kiện kết thúc, thuộc tính ngữ cảnh như scenario, protocol, service, symptom, tool và thống kê sử dụng như maturity, success count, failure count. Trạng thái bộ nhớ được lưu tại \code{runtime/memory/\{bank\_id\}/skills.json}.

Điểm mới của mô-đun là biểu diễn kinh nghiệm thành các quy trình chẩn đoán ngắn gọn, có điều kiện kích hoạt và có thể tái sử dụng. Khi bắt đầu một incident, hệ thống truy xuất kỹ năng phù hợp theo ngữ cảnh như scenario, protocol, service, symptom và tool. Điều này giúp tránh áp dụng sai kinh nghiệm, ví dụ không dùng quy trình kiểm tra BGP cho một case P4/BMv2 hoặc không dùng quy trình kiểm tra DHCP cho lỗi SDN controller.

Sau mỗi episode, hệ thống tính phần thưởng kết hợp từ detection, localization và RCA; lưu episode vào experience pool; sinh Semantic Gradient để rút kinh nghiệm; tạo nhiều ứng viên kỹ năng mới; rồi dùng PPO Gate phi tham số~\citep{ppo2017} để chỉ chấp nhận kỹ năng có cải thiện đủ rõ so với baseline. Điều kiện thành công nghiêm ngặt giúp giảm nguy cơ học vẹt đáp án benchmark hoặc ghi nhớ tri thức sai.

Cơ chế học được kiểm soát theo ba nguyên tắc. Thứ nhất, chỉ những episode đạt ngưỡng chất lượng mới được đưa vào golden pool, tránh để kết luận sai trở thành tri thức lâu dài. Thứ hai, kỹ năng mới phải vượt qua PPO Gate, tức là có lợi ích tương đối so với kỹ năng hiện có thay vì chỉ diễn đạt lại cùng một quy trình. Thứ ba, thư viện kỹ năng được giới hạn kích thước và có cơ chế nghỉ hưu kỹ năng kém hiệu quả, giúp bộ nhớ không phình to và không gây nhiễu ngữ cảnh cho LLM.

\subsection{Mô-đun Tool Refinement}

Tool Refinement tập trung vào việc giúp agent hiểu và sử dụng công cụ MCP tốt hơn, lấy cảm hứng từ DRAFT~\citep{draft2024}. Thay vì sinh công cụ mới hoặc sửa mã nguồn công cụ, hệ thống tự cải thiện phần mô tả công cụ dựa trên dữ liệu thực thi. Mỗi tài liệu công cụ gồm description, preconditions, hướng dẫn tham số, failure modes, usage notes, ví dụ đúng/sai, mastery score và convergence score. Trạng thái được lưu tại \code{runtime/tool\_evolution/\{library\_id\}/state.json}.

Quy trình học gồm ba bước. Explorer trích xuất các lần gọi công cụ từ \code{messages.jsonl}. Analyzer phát hiện khoảng trống hiểu biết như sai schema tham số, dùng tên thiết bị không tồn tại hoặc không biết diễn giải đầu ra. Rewriter viết lại tài liệu công cụ để bổ sung ràng buộc, ví dụ và hướng dẫn chẩn đoán. Khi mastery và convergence đạt ngưỡng, tài liệu được đóng băng adaptive để tránh viết lại không cần thiết.

Đề tài chọn cách tinh chỉnh tài liệu công cụ vì phù hợp với yêu cầu an toàn trong môi trường mạng. Công cụ nguyên thủy không bị thay đổi mã nguồn; agent chỉ học cách gọi đúng tham số, hiểu đúng tiền điều kiện và diễn giải đúng kết quả trả về. Điều này giảm rủi ro tạo ra thao tác không mong muốn trên thiết bị mạng, đồng thời vẫn cải thiện khả năng sử dụng tool theo thời gian.
